\section{Reorganizing the actual return before resetting the coupling}
\label{sec:physical-reset-reorganization}
We now turn the local return into the particular activity family used
in the scalar selection. This requires an identity on every field.
The difference between a local quadratic counterterm and the clipped
kinetic action has a contribution on unhealthy fields, which must
remain in the forced gas.

Fix the compact parameter interval, dictionary and local return of
Section~\ref{sec:selected-nonlinear-offering}. Work first at one
parameter value. Let $V=|\Lambda|$, let $0<t\le1$ be the coarse
threshold, and write $f_i$ for its regular entries, anchored at $o_i$
on supports $X_i$. Let $k_i$ be its forced entries. The budget has
radius $B\ge2$, support exponent $b>0$, derivative order $q\in\mathbb N$, $q\ge2$,
and spatial power $p\in\mathbb N$, $p\ge1$. The early regular majorants $u_i\ge0$ obey
\begin{equation}\label{eq:reset-early-majorant}
 \|f_i\|_{o_i,X_i}\le u_i,\qquad
 \sum_{i:o_i=o}e^{(b+3)c_B(X_i)}u_i\le H
                         \quad(o\in\Lambda),
\end{equation}
where $c_B$ is the covering number. The forced majorants obey the
rooted hit bound at exponent $b+1$, of cost $w$. All entries transform
under the same permutations of labels for translations and signed
lattice symmetries. The selected path has width $\sigma>0$, row and
higher-atom price $C$, and relaxation radius $r_c$ with
\[
             12\le r_c\sigma,\qquad (b+1)/B\le\sigma/4.
\]
Write $c_{\rm direct}$ for the direct coefficient of the clip. The
clipping parameters satisfy
\begin{equation}\label{eq:reset-clip-conditions}
 c_{\rm direct}\le1,\qquad 2\le t\,e^{(B-2)/r_c}.
\end{equation}
The first inequality is a consequence of the metric normalization
of the selected path recalled below. The second makes the kept-row
test local. Recall that the row at an anchor $o$ is kept when every
plaquette $q$ has chord at most
$t\exp(\operatorname{dist}(o,q)/r_c)$. Chords are at most $2$, so
\eqref{eq:reset-clip-conditions} makes this requirement automatic
for every plaquette at distance at least $B-2$ from $o$. The
remaining plaquettes have all four corners in the radius-$B$ halo of
$o$. Hence the test depends only on the field in that halo, and it
holds whenever the halo is healthy, because the chords read there are
then at most $t$. These are the conditions already used for the
actual clipped path. The torus side is greater than $20$.

\subsection{The scalar coefficient extracted from the local Hessians}
We first record how the Hessian of a regular entry acts on affine
fields. Let $f$ be a regular entry on a support $X$ which is
nonwrapping at its anchor $o$, and let $D^2f$ denote the second
derivative at zero of $A\mapsto f(\exp A)$, where $A$ runs over
Lie-algebra valued link fields. An affine field anchored at $o$ is
a link field
\[
 A_j(x)=a_j+\sum_k(x-o)_k\,b_{jk},
\]
with Lie-algebra coefficients $a_j,b_{jk}$; the difference $x-o$ is
read in $\mathbb Z^4$, which is unambiguous on a nonwrapping support.
Its curl is the constant two-form $w_{jk}=b_{kj}-b_{jk}$, $j<k$.
We claim that $D^2f(A,A')$, for two affine fields, depends only on
their curls. There are four steps.

First, $Df=0$ at the identity. A constant gauge transformation acts
on link directions by the adjoint action on every link. Invariance
of $f$ makes the restriction of $Df$ to one link an invariant linear
functional on the Lie algebra. Such a functional vanishes on
commutators, and a simple Lie algebra is spanned by its commutators.

Second, $D^2f(\nabla\phi,u)=0$ for every Lie-algebra valued site
function $\phi$ and every direction $u$, where
$(\nabla\phi)_j(x)=\phi(x+e_j)-\phi(x)$. Indeed, let $g_s$ be the
gauge transformation $x\mapsto\exp(-s\phi(x))$. Near zero the field
$g_s\cdot\exp A$ has a smooth logarithm $c(A,s)$ with $c(0,0)=0$ and
derivative $(A,s)\mapsto A+s\nabla\phi$ there. Gauge invariance
says $f(\exp c(A,s))=f(\exp A)$. Differentiate this identity twice
at zero. Since $Df=0$ at the identity, the second derivative of the
chart $c$ does not enter, and
\[
 D^2f(A+s\nabla\phi,A'+s'\nabla\phi)=D^2f(A,A')
\]
for all $(A,s)$ and $(A',s')$. Take $(A,s)=(0,1)$ and
$(A',s')=(u,0)$.

Third, $D^2f$ is symmetric, and $D^2f(A,A')$ is unchanged when $A$
or $A'$ is modified on links not contained in $X$, since $f$ reads
only the links of $X$.

Fourth, on a nonwrapping support a curl-free affine field is a
lattice gradient. If $b_{jk}=b_{kj}$, put
\[
 \phi(x)=\sum_j(x-o)_j\bigl(a_j-\tfrac12b_{jj}\bigr)
       +\tfrac12\sum_{j,k}(x-o)_j(x-o)_kb_{jk}.
\]
Then $\phi(x+e_d)-\phi(x)=a_d-\tfrac12b_{dd}
+\tfrac12\bigl(2\sum_k(x-o)_kb_{dk}+b_{dd}\bigr)=A_d(x)$ on every
link of the support.

Now let $A,\widetilde A$ be affine fields with the same curl. Their
difference is a curl-free affine field, hence agrees on the links of
$X$ with a gradient. The second and third steps give
$D^2f(A,u)=D^2f(\widetilde A,u)$ for every $u$, and symmetry gives
the same statement in the second argument. This proves the claim.

For a plaquette pair $a=(i,j)$, $i<j$, and a Lie-algebra element
$v$, let $W_a(v)$ be the affine field whose links in direction $j$
have the value $(x-o)_iv$ and whose other links vanish. Its curl is
$v$ in the plane $a$ and zero in the other planes. Every affine field
with curl $w$ therefore has the same curl as $\sum_aW_a(w_a)$.
The bilinear form $(v,v')\mapsto D^2f(W_a(v),W_d(v'))$ is invariant
under the adjoint action, by constant gauge invariance. Simplicity
of the Lie algebra identifies every
invariant bilinear form with a scalar multiple of the invariant
inner product. Hence there are real numbers $Q_f(a,d)$ with
\begin{equation}\label{eq:reset-affine-hessian}
 D^2f(A,A')=\sum_{a,d}Q_f(a,d)\langle w_a,w'_d\rangle
\end{equation}
for all affine fields $A,A'$ anchored at $o$, with curls $w,w'$.
Symmetry of the Hessian gives $Q_f(a,d)=Q_f(d,a)$.

For each $i$ put $Z_i=X_i\cup\operatorname{halo}_2\{o_i\}$.
On a nonwrapping $Z_i$, let $Q_i(a,d)=Q_{f_i}(a,d)$ be the
coefficient in \eqref{eq:reset-affine-hessian} for the entry $f_i$
and its anchor $o_i$. Set $Q_i=0$
when $Z_i$ wraps. This convention is preserved by the label symmetries.
The Hessian estimates give
\begin{equation}\label{eq:reset-projection-price}
 t^2|Q_i(a,a)|\le u_i,\qquad t^2|Q_i(a,d)|\le3u_i.
\end{equation}
Indeed, test a unit colour vector on one affine curvature direction,
scaled by $t$, for the first inequality. Its link size is bounded by
$t\operatorname{dist}(e_-,o)$, hence by
$t(1+\operatorname{dist}(e_-,o))^p$, so it is an admissible direction.
For the second use the
direction $t(A_a+A_d)/2$. Expanding its symmetric Hessian gives four
times this diagonal test minus the two single-direction tests, divided
by two. The derivative bounds give $(4+1+1)u_i/2=3u_i$.
For wrapping supports the inequalities hold because $Q_i=0$.

The series $Q_o(a,d)=\sum_{i:o_i=o}Q_i(a,d)$ converges absolutely
by \eqref{eq:reset-early-majorant}. Reindexing the label series by a
translation shows that $Q_o$ is independent of $o$. Reindexing by a
reflection of an axis lying in exactly one of two distinct plaquette
pairs changes the sign of the corresponding off-diagonal coefficient,
so that coefficient is zero. An unsigned permutation carries any
plaquette pair to any other, so all diagonal coefficients agree.
Consequently there is one real number $\Delta$ with
\begin{equation}\label{eq:reset-isotropic-coefficient}
                   Q_o(a,d)=\Delta\,\delta_{ad},
                       \qquad t^2|\Delta|\le H.
\end{equation}
The last bound uses the diagonal estimate in
\eqref{eq:reset-projection-price} at a single anchor; no volume factor
occurs.

The curvature read in the counterterm is the signed clover. For an
axis pair $a=(i,j)$, put
\[
 \phi_a(o,U)=\frac14\sum_{\epsilon,\eta\in\{\pm1\}}
 \epsilon\eta\,\operatorname{curv}_G
 \bigl(U(o;\epsilon e_i,\eta e_j,-\epsilon e_i,-\eta e_j)\bigr),
\]
where the argument is the holonomy of the indicated four-step loop
based at $o$, and $\operatorname{curv}_G$ is the fixed normalized
Lie-algebra projection. Each loop lies in the radius-two halo. A gauge
transformation acts by the adjoint at $o$; a signed lattice symmetry
permutes these four loops and their signs. This gives exactly the
signed two-form transformation law at the common anchor. On the
exponential line of an affine link field of constant curvature $w$,
the derivative of each oriented loop's curvature at zero is
$\epsilon\eta w_a$. The average therefore has derivative $w_a$,
with no extra normalization factor. Define
\[
 P_i(U)=\frac12\sum_{a,d}Q_i(a,d)
                  \langle\phi_a(o_i,U),\phi_d(o_i,U)\rangle,
 \qquad r_i=f_i-f_i(1)-P_i.
\]
The affine Hessian of $r_i$ vanishes when $Z_i$ is nonwrapping, and
$r_i(1)=0$. To see the first assertion, let $A,A'$ be affine fields
anchored at $o_i$, with curls $w,w'$. Each $\phi_a(o_i,\cdot)$
vanishes at the identity and has derivative $w_a$ along $A$, as just
shown. The product rule and $Q_i(a,d)=Q_i(d,a)$ give
\[
 D^2P_i(A,A')=\sum_{a,d}Q_i(a,d)\langle w_a,w'_d\rangle .
\]
By \eqref{eq:reset-affine-hessian} this equals $D^2f_i(A,A')$, and
the constant $f_i(1)$ has zero Hessian. Hence $D^2r_i(A,A')=0$.
No assertion is made for wrapping $Z_i$, where $Q_i=0$.

Write $\chi_X$ for the healthy-set indicator on $X$. A site is
unflagged precisely when every plaquette having it as a corner is
healthy. In particular $\chi_{\{o_i\}}=1$ makes the four clover
plaquettes at $o_i$ healthy.
Define the supported discrepancy
\[
 d_i=f_i(1)+\chi_{\{o_i\}}P_i-\chi_{X_i}f_i
                                  +\chi_{Z_i}r_i.
\]
It vanishes when $Z_i$ is healthy: all three indicators then equal
one, and the four terms cancel by the definition of $r_i$. When
$Z_i$ is not healthy the last term is absent. In both cases
\[
 \chi_{X_i}f_i=f_i(1)+\chi_{\{o_i\}}P_i-d_i+\chi_{Z_i}r_i .
\]
This has the form of the masked splitting of
Section~\ref{sec:local-splitting}, but it is not a special case of
that identity. There the extracted section is masked by the health
of the stencil on which it is supported. Here the quadratic
counterterm is masked by the health of the single anchor $o_i$,
which is what allows the counterterms of all entries with a common
anchor to be summed below. The displayed identity is therefore
verified directly, as was just done. Put
\[
 c_0=\sum_i f_i(1),\quad D=\sum_i d_i,\quad
 R_0=\sum_i\chi_{Z_i}r_i,\quad
 h(U)=\frac12\sum_o\chi_{\{o\}}(U)\sum_a|\phi_a(o,U)|^2.
\]
All sums converge absolutely. For example each summand is bounded by
a fixed multiple of $u_i$, and summing
\eqref{eq:reset-early-majorant} over the finite set of anchors bounds
$\sum_i u_i$ by $VH$. The isotropic coefficient identity now gives
\begin{equation}\label{eq:reset-first-split}
                \sum_i\chi_{X_i}f_i=c_0+\Delta h-D+R_0.
\end{equation}
The common singleton indicator at each anchor is essential to summing
the quadratic counterterms before using isotropy.

\subsection{Replacing the quadratic counterterm by the selected path}
Write $\mathcal K$ for the actual clipped kinetic action and
$\mathcal L$ for its loop part. The early return is
\begin{equation}\label{eq:reset-early-density}
 e^{V\tau-\beta\mathcal K(U)-\mathcal L(U)
             -\sum_i\chi_{X_i}(U)f_i(U)-\alpha\mathcal Q(U)}
                            \mathcal G(U).
\end{equation}
Here $\mathcal G$ is the early forced gas and $\mathcal Q$ is the
masked low-channel presentation of the quadratic action. It includes
the kept-row halo in each support. Thus it is specified before the
reset; it is not the full kinetic action on arbitrary fields.

We describe the needed identity channel by channel. A low channel
is either a plaquette term or an ordered pair of entries in one row
of the selected path. Denote its coefficient by $a_d$, its unmasked
quadratic field by $P_d$, its local constant-curvature substitute by
$Q_d$, and its actual clipped or kept field by $L_d$. For a row pair,
\[
 a_d=c_i c_j,\qquad
 P_d=\tfrac12\langle\phi_i,\phi_j\rangle,
 \qquad Q_d=\tfrac12\langle\phi_{a_i},\phi_{a_j}\rangle.
\]
The latter curvature reads are at the common anchor. For a plaquette,
$a_d$ is the direct coefficient, $P_d$ is its plaquette action and
$Q_d$ is half the squared local curvature. Let $Y_d$ be the union of
the radius-$B$ anchor halo with the two walks when present. On its
healthy set, $L_d=a_dP_d$. Define
\begin{align*}
 F_d&=\Delta a_d(Q_d-P_d),\\
 J_d&=(\alpha+\Delta)L_d
              -\Delta a_d\chi_{\{o\}}Q_d
                    \quad\hbox{on the complement of the healthy set of }Y_d,
\end{align*}
and set $J_d=0$ on that healthy set. Directly in these two cases,
\begin{equation}\label{eq:reset-low-identity}
 (\alpha+\Delta)L_d+\chi_{Y_d}F_d
   =\alpha\chi_{Y_d}a_dP_d
                      +\Delta a_d\chi_{\{o\}}Q_d+J_d.
\end{equation}
This also proves that the discrepancy is supported on an unhealthy
$Y_d$. No derivative of its indicator is needed.

Average these identities over signed symmetries and sum over
translations, with normalization $1/|\mathcal S|$ for the signed group
$\mathcal S$. This normalization is applied to every channel, including
its discrepancy. To identify the quadratic substitute, let $S(0)$ be the actual
zero-momentum row matrix. Absolute row summability permits the
row-pair sum to be regrouped as
\[
 \sum_d a_d Q_d(w)
   =\tfrac12\{c_{\rm direct}\|w\|^2+\|S(0)w\|^2\}
   =\tfrac12\|w\|^2.
\]
The second equality is the selected metric normalization
$c_{\rm direct}I+S(0)^TS(0)=I$. Multiplying by the singleton healthy
indicator and summing anchors gives $h$; the signed average preserves
the squared norm of the clover. Thus the averaged sum of
$a_d\chi_{\{o\}}Q_d$ is exactly $h$. The averaged sum of
$\chi_{Y_d}a_dP_d$ is $\mathcal Q$.
For every cubic or quartic path atom $A_d$, add the regular correction
$-(\alpha+\Delta)A_d$ on its original mask. These atoms have zero
affine Hessian because their curvature factors vanish at the identity
and their degrees are at least three. Denote the complete averaged
regular correction by $F$, and the averaged sum of $J_d$ by $J$.
Since the low and high channels sum to $\mathcal K$, equation
\eqref{eq:reset-low-identity} becomes
\begin{equation}\label{eq:reset-path-identity}
       (\alpha+\Delta)\mathcal K+F
                              =\alpha\mathcal Q+\Delta h+J.
\end{equation}
The low correction also has zero affine Hessian: on an affine field
at the identity, transport along each path has the same first
curvature derivative as its constant-curvature substitute. Thus
$Q_d-P_d$ has zero second derivative there. Polarization gives the
mixed assertion. This argument is only used on nonwrapping supports,
where the affine field and its walks have their specified lifts.

Combining \eqref{eq:reset-first-split} and
\eqref{eq:reset-path-identity} in \eqref{eq:reset-early-density} gives
\begin{equation}\label{eq:reset-complete-density}
 e^{V(\tau-c_0/V)-(\beta+\alpha+\Delta)\mathcal K
                          -\mathcal L-(R_0+F)}
                                     \mathcal G\,e^{D+J}.
\end{equation}
Every sign follows by substituting the two displayed identities.
This equality holds even where $\mathcal G$ vanishes or changes sign.
Absorb $D+J$ into that gas using the supported exponential construction
of Section~\ref{sec:soft-log-fusion}. Its hypotheses hold because each
discrepancy vanishes on its healthy support. This constructs the actual
output at coupling $\beta+\alpha+\Delta$ and vacuum parameter
$\tau-c_0/V$.

\subsection{The regular and forced costs of this reorganization}
We give the numerical bounds so that this identity is accompanied
by estimates for its actual arrays. Let $C_{\rm pair}$ be the
fixed ordered-derivative price for a pair of local curvatures, and put
\[
 C_{\rm sub}=2+\tfrac32|\mathcal A|^2C_{\rm pair},
                  \qquad C_{\rm reset}=e^{81b}C_{\rm sub},
\]
where $\mathcal A$ is the six-element set of plaquette axes.
Equation~\eqref{eq:reset-projection-price} and the product rule bound
both the regular norm of $r_i$ and the value of $d_i$ by
$C_{\rm sub}u_i$. The added radius-two halo increases the cover number
by at most $81$, whence the regular cost is at most
$C_{\rm reset}H$. Its hit cost, denoted $S(H)$, is the footprint
payment with population radius $4B+2$, support increment $81$,
input exponent $b+3$ and output exponent $b+1$. In the notation of
the explicit footprint price $\mathcal P$ of
Section~\ref{subsec:ex-lattice-rg-crowd},
\[
 S(H)=\mathcal P(1,4B+2,0,1,C_{\rm sub},H,b+1,81).
\]
This is the ordinary application, at exponent $b+1$ and width
reserve one; it is not the raised application displayed in that
section. It is linear
in $H$ and has only polynomial dependence on $B$ at fixed stage.

Let $C_{\rm turn}$ be the ordered-derivative cost of a signed
pullback. For the low and high channels let $L_0,L_1$ be their rooted
regular sums after support weighting but before this pullback cost.
Let $L_h$ be the low-channel hit sum, with the relaxation factor
$e^{\ell_d/r_c}$ included. The channel estimates give
\begin{align}
 \|F_{\rm low}\|&\le C_{\rm turn}|\Delta|t^2L_0,\nonumber\\
 \|F_{\rm high}\|&\le C_{\rm turn}|\alpha+\Delta|t^3L_1,
                                                    \label{eq:reset-channel-costs}\\
 \|J\|_{\rm hit}&\le
 \left(2t|\alpha+\Delta|+\tfrac12C_{\rm pair}t^2|\Delta|\right)L_h.
                                                        \nonumber
\end{align}
Here is the summability behind these costs. A low support satisfies
$Bc_B(Y_d)\le85B+\ell_d$ and a high support satisfies
$Bc_B(Y_d)\le9B+\ell_d$. The weighted row-pair sum is the product
of the two weighted row sums, hence is bounded by $C^2$; the direct
plaquette contribution adds its fixed coefficient. Derivatives of
fixed order multiply these charges by a fixed polynomial in
$1+\ell_d$. The inequality
$(1+\ell)^k\le C_{k,\epsilon}e^{\epsilon\ell}$, obtained by
maximizing $(1+x)^ke^{-\epsilon x}$ on $x\ge0$, pays this polynomial
from half the width. The remaining half pays the support weight.
For the hit sum, $r_c^{-1}\le\sigma/12$ and
$(b+1)/B\le\sigma/4$ leave positive width for the rooted footprint
sum. These statements hold term by term with absolute coefficients,
so also for the countable channel sums.

The low regular estimate in \eqref{eq:reset-channel-costs} follows
from the two curvature factors in $P_d,Q_d$. For its discrepancy,
the kept-row bound is
$|L_d|\le2t e^{\ell_d/r_c}|a_d|$, and the singleton-masked
substitute has size at most $C_{\rm pair}t^2/2$.
The triangle inequality gives the third estimate. The high estimate
uses $t^4\le t^3$ and the cubic and quartic path bounds.
Finally \eqref{eq:reset-isotropic-coefficient} and $t\le1$ imply
\[
 |\alpha+\Delta|t^3\le |\alpha|t^3+H,
 \qquad
 2t|\alpha+\Delta|+\tfrac12C_{\rm pair}t^2|\Delta|
 \le2t(|\alpha|+H/t^2)+\tfrac12C_{\rm pair}H.
\]
Thus the full regular and soft costs are bounded respectively by
\begin{align*}
 E_{\rm new}&=C_{\rm reset}H+
       C_{\rm turn}\{HL_0+(|\alpha|t^3+H)L_1\},\\
 S_{\rm new}&=S(H)+
       \{2t(|\alpha|+H/t^2)+\tfrac12C_{\rm pair}H\}L_h.
\end{align*}
The reorganization requires two numerical purchases for its output
allowances $\delta$ and $\gamma$. The regular one is
$E_{\rm new}\le\delta$. The forced one asks for numbers
$\mu,\kappa\ge0$ with
\begin{equation}\label{eq:reset-forced-purchase}
 S_{\rm new}\le\mu\le1,\qquad
 (2B+1)^4I\le1,\qquad
 e^{-\kappa}(e^{I}-1)\le\gamma,\qquad I=e^\kappa w+2\mu .
\end{equation}
These are the hard/soft absorption inequalities
\eqref{eq:nonlinear-payment} for the pair $(w,\mu)$ and target
$\gamma$. Under them the output has
regular cost $\delta$ and forced cost $\gamma$, with the original
locality, symmetry and affine-Hessian conditions. These numerical
purchases must be imposed on this same output budget.

\subsection{The stage constants and the early allocation}
The regular purchase is obtained by fixing in advance how small the
early price $H$ must be. The coefficients
$C_{\rm reset},C_{\rm turn},L_0,L_1$ depend only on the derivative
order, spatial power, support exponent, clip, width and path price
of the output stage. They do not depend on the radius $B$ or the
threshold $t$: the cover-number increments $81$, $85$, $9$ and the
polynomial derivative factors above were bounded without reference
to either. Put
\begin{equation}\label{eq:reset-stage-constants}
 B^{\rm reset}=C_{\rm reset}+C_{\rm turn}(L_0+L_1),\qquad
 v^{\rm reset}=C_{\rm turn}L_1,\qquad
 c^{\rm reset}=\frac1{64(1+B^{\rm reset})}.
\end{equation}
Collecting the coefficients of $H$ and of $|\alpha|t^3$ in the
display for $E_{\rm new}$ gives the identity
\[
 E_{\rm new}=B^{\rm reset}H+v^{\rm reset}|\alpha|t^3 .
\]
Since $B^{\rm reset}\ge0$,
\begin{equation}\label{eq:reset-credit-bounds}
 0<c^{\rm reset}\le\frac1{64},\qquad
 B^{\rm reset}c^{\rm reset}\le\frac1{64}.
\end{equation}
Now let $\delta_E>0$ be the regular allowance of the output level,
and let $A\ge0$. Suppose
\begin{equation}\label{eq:reset-allocation}
 H=c^{\rm reset}\delta_E,\qquad |\alpha|\le A,\qquad
 64\,v^{\rm reset}A\,t^3\le\delta_E .
\end{equation}
Then \eqref{eq:reset-credit-bounds} gives
\[
 E_{\rm new}\le\frac{\delta_E}{64}+\frac{\delta_E}{64}
           =\frac{\delta_E}{32}.
\]
Thus the regular purchase holds with $\delta=\delta_E/32$. The first
condition in \eqref{eq:reset-allocation} is an allocation: it is the
fraction of the output allowance which the local return of
Section~\ref{sec:selected-nonlinear-offering} is asked to meet, and
is the constant called $c_{\rm st}$ there. It is a choice made
before the block length, not a consequence of the reorganization.
The third condition is a stage-entry purchase. With the uniform
bound $A=A_j$ of the next subsection and
$t=x^{-1/2}e^{(\log\log x)^2}$, $\delta_E=x^{-11/10}$, the ratio
$t^3/\delta_E=x^{-2/5}e^{3(\log\log x)^2}$ tends to zero, so the
condition holds beyond an entry threshold depending on the stage
and on the block length through $A_j$.

Section~\ref{sec:nonlinear-payment-selection} proves the soft-price
estimate \eqref{eq:nonlinear-soft-price}, which bounds $S_{\rm new}$
by a stage constant times $(1+B)^4t$ when $|\alpha|\le A$ and
$H\le t^2$, and the elementary payment \eqref{eq:nonlinear-payment}.
It does not contain the regular inequality $E_{\rm new}\le\delta$;
that is the allocation just described. The comparison of
$H,\delta,\gamma$ with the prices of the selected recipe is proved in
Section~\ref{sec:common-funded-history}. The instantiation of
\eqref{eq:reset-forced-purchase} at the particular forecast pair
$(w,\gamma)$ of the selected history remains an input here: it
requires witnesses $\mu,\kappa$ satisfying all its inequalities for
that pair. Section~\ref{sec:common-funded-history} describes this
purchase in words but does not exhibit the witnesses. The present
section is conditional on \eqref{eq:reset-allocation} and
\eqref{eq:reset-forced-purchase}.

\subsection{Continuity in the parameter}
Let the parameter $y$ range over the compact interval $T$ fixed at
the start, and write $f_i(y),k_i(y),\Delta(y)$ for the early arrays
and the extracted coefficient. The output arrays must be continuous
in $y$ in the labelled regular and forced norms. This does not follow
from continuity of finitely many operations, because the forced
absorption involves infinitely many labels. We use the two
comparison properties of the early arrays supplied by
Section~\ref{sec:selected-nonlinear-offering}, both in the one
dictionary fixed there:
\begin{itemize}
\item for every $y\in T$ and every label, the regular entry
 $f_i(y')$ converges to $f_i(y)$ in its regular norm, and the forced
 entry $k_i(y')$ converges to $k_i(y)$ uniformly in the field, as
 $y'\to y$ in $T$;
\item every sequence in $T$ converging to $y$ has a subsequence along
 which the regular norms of the $f_i$ and the suprema of the $k_i$
 have summable majorants $e_i,h_i$ common to the subsequence and its
 limit.
\end{itemize}
On the fixed finite torus every support has cover number at most the
number of sites, so the support weights are bounded and plain
summability of majorants is all that is needed below.
The first step is the absorption itself.

\begin{lemma}[Continuity of the forced absorption at a fixed period]
\label{lem:reset-absorption-continuity}
Fix a finite torus, a radius $B$, a support exponent $b$, and one
dictionary of hard labels $i$ and logarithmic labels $j$ as in
Section~\ref{sec:soft-log-fusion}, with supports, inventories and
eligibility tests independent of a parameter $y\in T\subseteq
\mathbb R$. Let $K_i(y,\cdot)$ and $\ell_j(y,\cdot)$ be the hard
factors and logarithms. Assume:
\begin{enumerate}
\item for every label and every $y\in T$,
 $\sup_U|K_i(y',U)-K_i(y,U)|\to0$ and
 $\sup_U|\ell_j(y',U)-\ell_j(y,U)|\to0$ as $y'\to y$ in $T$;
\item there are numbers $p_i,v_j\ge0$, independent of $y$, with
 $|K_i(y,U)|\le p_i$ and $|\ell_j(y,U)|\le v_j$ for all $y\in T$ and
 all fields $U$, and with $\sum_ip_i+\sum_jv_j<\infty$.
\end{enumerate}
Then the absorbed activities $\widehat K_C(y,\cdot)$ of
\eqref{eq:fusion-new-activity} satisfy
\[
 \sum_Ce^{b\,c_B(Z_C)}\sup_U
 \bigl|\widehat K_C(y',U)-\widehat K_C(y,U)\bigr|\longrightarrow0
 \qquad(y'\to y\ \hbox{in }T).
\]
\end{lemma}
\begin{proof}
At every field, $\widehat K_C(y,U)$ is the product of two indicators
that do not depend on $y$ and of the finitely many factors
$r_v(y,U)$, $v\in C$, where $r_i=K_i$ and $r_j=e^{\ell_j}-1$.
Put $a_i=p_i$ and $a_j=e^{v_j}-1$. Each factor is bounded by $a_v$
uniformly in $y$ and $U$, because $|e^x-1|\le e^{|x|}-1$, and by the
first hypothesis each factor converges uniformly in $U$. Writing the
difference of two finite products as a telescoping sum shows that
each term of the series tends to zero. Each term is at most
$2e^{b|\Lambda|}\prod_{v\in C}a_v$, since a cover number is at most
the number $|\Lambda|$ of sites. The numbers $a_v$ are summable:
$a_j\le v_je^{v_j}$ and the $v_j$ are bounded. The labels $C$ are
finite sets of constituent labels, so
\[
 \sum_C\prod_{v\in C}a_v\le\prod_v(1+a_v)-1
                       \le\exp\Bigl(\sum_va_v\Bigr)-1<\infty .
\]
Dominated convergence for series gives the claim.
\end{proof}

The bound obtained in this proof depends on the number of sites. It
is used only for continuity at a fixed period, never as a
volume-independent price. The volume-independent forced cost of the
output is the separate estimate \eqref{eq:fusion-site-output}, paid
by \eqref{eq:reset-forced-purchase}.

We apply the lemma as follows. First, $\Delta$ is continuous on $T$.
The difference of two early arrays in the same dictionary is again a
family of regular entries, and the projection is linear. Applying
\eqref{eq:reset-isotropic-coefficient} to the difference gives
$t^2|\Delta(y')-\Delta(y)|\le d$ whenever the regular norms of the
differences $f_i(y')-f_i(y)$ have majorants whose support-weighted
sum over all labels is at most $d$. Such a $d$ can be taken
arbitrarily small for $y'$ near $y$: otherwise there would be a
sequence $y_n\to y$ along which this fails, and on its dominated
subsequence the entrywise convergence and the common summable
majorant $2e_i$ would give a contradiction by dominated convergence.

Second, every output entry depends continuously on $y$. The entries
$r_i$ and $d_i$ are linear in $f_i$, with norm at most $C_{\rm sub}$
times that of $f_i$, so they inherit the entrywise convergence. The
path corrections $F_d$ and $J_d$ are fixed functions multiplied by
$\Delta$ or by $\alpha+\Delta$, so they are continuous because
$\Delta$ is.

Third, let $y_n\to y$ and pass to the subsequence with common early
majorants $e_i,h_i$. Along it $|\Delta(y_n)|\le|\Delta(y)|+1$
eventually. The output regular and logarithmic entries then have the
common summable majorants $C_{\rm sub}e_i$ for the subtracted entries
and the channel bounds of \eqref{eq:reset-channel-costs} with
$|\Delta|$ replaced by $|\Delta(y)|+1$ for the path entries; the hard
entries keep the majorants $h_i$. Summability of the channel
bounds is the summability established after
\eqref{eq:reset-channel-costs}, which did not use the size of
$\Delta$.

Fourth, on this subsequence together with its limit point the
hypotheses of Lemma~\ref{lem:reset-absorption-continuity} hold: the
first by the second step, the second by the third step. The regular output
entries are not changed by the absorption, and their weighted
difference sums tend to zero by the same dominated convergence.
Hence the full output arrays converge along the subsequence. Since
every sequence converging to $y$ has such a subsequence, the output
arrays are continuous at $y$: a sequence violating continuity would
have a dominated subsequence still violating it.

\subsection{The two healthy tests and the actual reset}
Suppose the early cost obeys $H\le F/4$, where $F$ is the scalar
fund used in Section~\ref{sec:coupling-reset-selection}. On a globally
healthy field all masks are one and the early gas is one. The early
density is consequently
\[
 H_x(U)=\exp\{V\tau-\beta\mathcal K(U)
                         -\alpha\mathcal K_2(U)-\mathcal L(U)-E(U)\},
 \qquad |E(U)|\le VH.
\]
Here $\mathcal K_2$ is the quadratic part and $E$ the early regular
sum. At $U=1$, the three path terms vanish, so positivity is explicit
and
\[
 \left|V^{-1}\log H_x(1)-\tau\right|\le H,\qquad
 |E(U)-E(1)|\le2VH\le VF/2.
\]
The difference of the full kinetic action and its quadratic part has
size at most $2VCt^3$: each cubic and quartic atom is bounded by its
row charge times $t^3$ or $t^4$, respectively, and $t^4\le t^3$.
The selected path construction bounds
$|\alpha|\le A_j$, uniformly in cutoff and period, where
$A_j$ is its drift bound plus $|s_G\log L_j|$ for the fixed stage.
Purchase $2A_jCt^3\le F/2$. This is possible at each fixed stage since $t=x^{-1/2}e^{(\log\log x)^2}$, while $F$ is a
positive stage constant times $x^{-11/10}$; their ratio tends to
zero. Adding the two errors proves the precise surplus bound $VF$
at the nontrivial healthy test. The vacuum error is bounded by $F/4$.

The output arrays in \eqref{eq:reset-complete-density} depend on the
computed $\Delta$, and are fixed while the trial dial varies. Their
dial-motion cost is therefore zero. Their continuity in the original
parameter was proved in the preceding subsection. The scalar selection
therefore applies to these arrays. Its residual is
$(d-\Delta)\mathcal K(U_*)$, so the selected dial is exactly
$\Delta$. Its vacuum value is exactly $\tau-c_0/V$ by
\eqref{eq:reset-complete-density} at the identity.

Finally retain this constructed family as the witness for the
history step. The recipe conditions refer to its actual supports,
normalized density identity, regular and forced bounds, and bound
$t^2|\Delta|\le H$. They are supplied by the construction above and
the separately purchased recipe permissions. Equality of densities
by itself would not supply those permissions. Once the common funding
provides them, this is the continuous one-step map used in the
surviving-interval induction.
